\documentclass[11pt,reqno,a4paper]{amsart}
\usepackage[T1]{fontenc}
\usepackage{amssymb,booktabs}
\usepackage[hidelinks]{hyperref}
\allowdisplaybreaks[2]
\raggedbottom
\setlength{\emergencystretch}{1.8em}

\newtheorem{theorem}{Theorem}[section]
\newtheorem{proposition}[theorem]{Proposition}
\newtheorem{corollary}[theorem]{Corollary}
\newtheorem{lemma}[theorem]{Lemma}
\theoremstyle{remark}
\newtheorem{remark}[theorem]{Remark}
\newtheorem*{theoremT}{Theorem T}

\newcommand{\Fac}{\operatorname{Fac}}
\newcommand{\e}{\varepsilon}
\newcommand{\Y}{\mathcal Y}
\newcommand{\tern}[1]{\mathtt{#1}}
\newcommand{\R}{\mathbb R}
\newcommand{\Q}{\mathbb Q}
\newcommand{\Z}{\mathbb Z}
\newcommand{\dist}[1]{\lVert#1\rVert}
\newcommand{\rev}{\mathrm{rev}}
\newcommand{\word}[1]{\mathtt{#1}}
\newcommand{\HH}{\mathcal H}

\title[Affine separation and a transcendental family]
{Affine separation of squares and palindromes and a full-dimensional transcendental family}
\newsavebox{\affilbox}
\address{School of Artificial Intelligence, Shanghai Jiao Tong University, Shanghai, China}
\email{mqyqwerty@sjtu.edu.cn}
\date{}
\subjclass[2020]{11J68, 11J83, 68R15}
\keywords{Algebraic irrationals, integer-base expansions, squares, palindromes, rational affine maps, Hausdorff dimension, transcendence}
\hypersetup{pdfauthor={Qiyuan Ma},pdftitle={Affine separation of squares and palindromes and a full-dimensional transcendental family},pdfdisplaydoctitle=true}

\begin{document}
\sbox{\affilbox}{%
  \begin{minipage}{0.92\textwidth}
    \centering\normalfont\normalsize
    School of Artificial Intelligence,\\
    Shanghai Jiao Tong University, Shanghai, China\\
    ORCID: \href{https://orcid.org/0009-0000-9532-0821}{0009-0000-9532-0821}\\
    \href{mailto:mqyqwerty@sjtu.edu.cn}{mqyqwerty@sjtu.edu.cn}
  \end{minipage}}
\author[Qiyuan Ma]{Qiyuan Ma\\[6pt]\protect\usebox{\affilbox}}

\begin{abstract}
Bugeaud's Problem~10.42 asks whether every real algebraic irrational has unbounded square-root lengths and unbounded palindrome lengths in every integer-base expansion. We show that these two properties can be separated by rational affine maps. There are uncountably many irrationals whose rational affine images all have unbounded palindromes of both parities, while excluding long high-power repetitions; and there are uncountably many whose rational affine images all have unbounded square-root lengths in every grouped base, yet have bounded palindrome lengths. In particular, an argument invariant under rational affine maps cannot establish both conclusions at once. We also construct a set of transcendental numbers of Hausdorff dimension one such that, for every nonconstant rational function $f$ with rational coefficients, the value $f(x)$ is defined and, in every integer base, has bounded square-root lengths and bounded palindrome lengths. Thus any positive answer to the unbounded-length questions must use algebraicity in an essential way: the two properties fail on a full-dimensional set of real numbers. Both questions of Problem~10.42 remain open for algebraic irrationals.
\end{abstract}
\maketitle

\section{Introduction}\label{sec:intro}

For an integer $b\ge2$ and a real number $x$, write $\{x\}=x-\lfloor x\rfloor$ and $\dist{x}=\min_{n\in\Z}|x-n|$. Let $\mathsf S_b(x)$ mean that the canonical base-$b$ expansion of the fractional part of $x$ contains factors $VV$ with $|V|$ unbounded. Let $\mathsf P_b(x)$ mean that it contains palindromes of unbounded length. Bugeaud's Problem~10.42 asks whether each property holds for every real algebraic irrational $x$ and every integer $b\ge2$ \cite[p.~220]{Bugeaud}. Adamczewski and Rampersad posed the corresponding questions in base two \cite{AR}.

The digit block starting after $i$ places and of length $d$ has integer value
\begin{equation}\label{eq:digitblock}
 D_{i,d}(x)=\lfloor b^{i+d}x\rfloor-b^d\lfloor b^i x\rfloor.
\end{equation}
Leading zeros are included. For a word $W$ let $[W]_b$ denote its integer value. Deleting any fixed prefix preserves both $\mathsf S_b$ and $\mathsf P_b$. A long palindrome can be trimmed symmetrically. If all unbounded squares began in a finite set of starts, one start would have unbounded roots. The second half of a sufficiently large square at that start contains each earlier given square at the same start, and moves that square beyond the deleted segment.

The long-repetition criterion of Adamczewski--Bugeaud--Luca, as stated by Bilu \cite[Theorem~3.3, p.~8]{Bilu}, implies that the greatest length of two disjoint equal factors among the first $N$ digits of a fixed real algebraic irrational is $o(N)$. The corresponding length $L_b(N)$ with distinct starts and possible overlap also satisfies
\begin{equation}\label{eq:known-early}
 L_b(N)=o(N).
\end{equation}
Indeed, overlapping matches of length $\ell$ at distance $d<\ell$ give a period-$d$ segment of length $\ell+d$. It contains $\lfloor\ell/d\rfloor+1$ adjacent copies of a length-$d$ word. Splitting them into two equal groups gives a disjoint match of length at least $\ell/3$, as in \cite[Lemma~3.4]{Bilu}.

The two properties $\mathsf S_b$ and $\mathsf P_b$ are not interchangeable, and neither is an automatic feature of ``most'' real numbers. The purpose of this paper is to make both observations precise.

Theorem~\ref{thm:affine} separates the two unbounded-length assertions under rational affine maps. For any finite set of nonzero rational slopes, there is an uncountable family of irrationals $x$ such that every image $\lambda x+t$ has unbounded palindromes of both parities, while every rational affine image of $x$ excludes arbitrarily long high powers; and there is a complementary uncountable family for which every such image has unbounded square-root lengths in every grouped base $b^s$, yet has bounded palindrome lengths in each of those bases. Each family contains uncountably many transcendental numbers. Consequently, a proof of both $\mathsf S_b$ and $\mathsf P_b$ for algebraic irrationals cannot proceed by an argument that is invariant under rational affine maps and makes no further use of algebraicity.

Theorem~\ref{thm:nonlinear} supplies a complementary metric obstruction. There is a set $\mathcal K$ of Hausdorff dimension one, consisting entirely of transcendental numbers, such that for every $x\in\mathcal K$ and every nonconstant $f\in\Q(X)$, the value $f(x)$ is defined and its expansion in every integer base has bounded square-root lengths and bounded palindrome lengths. In particular, the unbounded-length conclusions of Problem~10.42 fail throughout a full-dimensional set, and they fail not only for $x$ itself but for every nonconstant rational-function image of $x$. A positive answer for algebraic irrationals must therefore use algebraicity in an essential way. The corresponding polynomial constraints in the construction exclude algebraic points from this particular family; it remains unknown whether the original avoidance sets contain any algebraic irrational.

The algebraic side of the comparison is represented by a congruence restriction that every algebraic irrational must satisfy. Corollary~\ref{cor:modular}, deduced from a cofactor form of the Subspace Theorem, asserts that early length-$\ell$ blocks of a fixed algebraic irrational are pairwise distinct modulo a large prime power. Proposition~\ref{prop:modtest} exhibits, within a single bounded-avoidance class, both uncountable families that obey this restriction and a positive-dimensional compact class that violates it at unbounded lengths and is therefore transcendental.

Theorem~\ref{thm:moving-maps} extends the avoidance from a fixed countable family of maps to moving families of controlled cardinality, still at dimension arbitrarily close to one. The supporting cylinder-event lemmas are recorded in Section~\ref{sec:sparse-arithmetic}.

Both unbounded-length questions of Problem~10.42 remain open for every algebraic irrational in every integer base.

\section{An early congruence criterion}\label{sec:arith}

The following cofactor bound will be used to test constructions against an arithmetic constraint that every algebraic irrational must satisfy.
For a finite set of primes $S$, put $[E]_S=\prod_{p\in S}p^{v_p(E)}$ and $w_S(E)=|E|/[E]_S$ for $E\ne0$. Set $w_S(0)=1$.

\begin{theorem}\label{thm:cofactor}
Fix $b$, a finite prime set $S$ containing the prime divisors of $b$, real algebraic $\beta_1,\ldots,\beta_t$, a finite number of error terms, and constants $A,C>0$. For every $0<\delta<1$ the following holds at sufficiently large $\ell$. If integers $U_i,E_j,P$ satisfy $|U_i|,|E_j|\le Cb^{A\ell}$,
\[
 \Phi=\sum_i\beta_i U_i\notin\Q,\qquad
 \left|b^\ell(\Phi-P)-\sum_jE_j\right|\le C,
\]
then
\[
 \prod_i w_S(U_i)\prod_j w_S(E_j)>b^{(1-\delta)\ell}.
\]
\end{theorem}

\begin{proof}
Suppose to the contrary that there is an infinite sequence of increasing $\ell$. Write $X_i=b^\ell U_i$, $Z=b^\ell P$, and $L=\sum_i\beta_iX_i-Z-\sum_jE_j$. Pass to a subsequence with fixed zero pattern, remove identically zero coordinates, and keep $|L|\le C$. All coordinates have height $O(b^{(A+1)\ell})$, including $Z$ by the assumed error bound.

At each stage the value $b^\ell\Phi$ is a linear combination of the retained $X_i$, with fixed algebraic coefficients. For $L$ the coefficients of these coordinates differ from those of $b^\ell\Phi$ by rational numbers. All coefficients of $Z,E_j$ in $L$ are rational. In particular some retained $X_i$ has irrational coefficient in $L$. Drop $Z$ if its coefficient is zero.

Use all coordinate forms at each $p\in S$. At the real place replace $Z$ by $L$ if $Z$ is present, and replace an $X_i$ with irrational coefficient otherwise. The forms are independent at every place. The coordinate replaced at the real place is a nonzero multiple of $b^\ell$, so its finite-place factor is at most $b^{-\ell}$. Each other original $X_i$ contributes $w_S(U_i)$ and each other error coordinate contributes $w_S(E_j)$. The total local product is at most $Cb^{-\delta\ell}$. Applying Schlickewei's Subspace Theorem in the form stated in \cite[Theorem A]{ABind} produces a fixed nontrivial rational relation along an infinite subsequence.

Eliminate an $E_j$ if any has nonzero coefficient in that relation. Then eliminate $Z$ if it has nonzero coefficient and the relation contains no errors. Either operation alters only the rational part of $L$. In the remaining case eliminate an $X_i$ using a relation among the $X_i$ and perform the same elimination in $b^\ell\Phi$. It preserves the asserted coefficient differences. All remaining coordinates are original coordinates and removing a cofactor can only decrease the product budget. The process strictly decreases the dimension. In dimension one the remaining coordinate is an $X_i$ with nonzero coefficient and magnitude at least $b^\ell$, contradicting $|L|\le C$.
\end{proof}

\begin{corollary}\label{cor:modular}
Fix a real algebraic irrational $\alpha$, an integer base $b$, a prime $p$, and $A,c>0$. For all sufficiently large $\ell$, the integer block values $D_{n,\ell}(\alpha)$ for $0\le n\le A\ell$ are pairwise distinct modulo $p^{\lceil c\ell\rceil}$.
\end{corollary}

\begin{proof}
For $n<m$, put $E=D_{m,\ell}(\alpha)-D_{n,\ell}(\alpha)$, $U_1=b^m$, $U_2=b^n$, and $P=\lfloor b^m\alpha\rfloor-\lfloor b^n\alpha\rfloor$. The error is at most $1$, and $\Phi=(b^m-b^n)\alpha$ is irrational. Theorem~\ref{thm:cofactor} gives $E\ne0$ and $w_S(E)>b^{(1-\delta)\ell}$. Since $|E|<b^\ell$, $[E]_S<b^{\delta\ell}$. Choose $p\in S$ and $\delta<c\log_b p$.
\end{proof}

The case $E=0$ in the corollary is already excluded by \eqref{eq:known-early}, since both blocks lie among the first $\lfloor (A+1)\ell\rfloor$ digits. The additional content of the corollary is a divisibility constraint for $E\ne0$, including when $p$ does not divide the base.

\section{Bounded avoidance and affine separation}\label{sec:construct}

\subsection{Avoiding examples and an exact algebraic identity}
Several supplemental constructions are useful for testing the quantifiers. In binary, the separate minimum caps are root length $2$ and palindrome length $5$. Dassow, Manea, Merca\c{s}, and M\"uller already obtained both caps in the same infinite binary word, and the caps $1,2$ in ternary \cite[Theorems 4 and 5]{DMMM}. A different synchronized ternary-to-binary code in the supplement attains the binary caps and has independent alternatives of fixed widths $605$, $619$, and $630$. It uses the Brinkhuis choices of Ekhad and Zeilberger with their correction \cite{EZ}. The finite verification checks the choice lengths and synchronization. The infinite proof transfers a long synchronized square to a square in a square-free source. A separate absolute-game construction works simultaneously for all integer bases and all fixed nonconstant algebraic-coefficient rational functions on their domains. It yields bounded palindrome and square lengths compatible with ordinary irrationality exponent $2$ and the stated restricted-denominator and repetition conditions. The supporting results use the potential-game equivalence and countable intersection properties \cite{FSU,BNY}.

The digit series also has an exact arithmetic constraint. Write $F(z)=\sum_{i\ge1}a_i z^i$ for the fractional digits in a base $b$. For an integer polynomial $P$, let $c_i=[z^i]P(F(z))$. If $P(F(1/b))=0$, then
\[
 R_m=\sum_{i>m}c_i b^{m-i}=-\sum_{i\le m}c_i b^{m-i}\in\Z,
 \qquad b R_{m-1}=c_m+R_m .
\]
This is the tail integrality used in \cite{BBCP}. In the supplemental constructions with isolated convolution coefficients, two long zero buffers and a nonzero intervening block put one such integer strictly between consecutive integers. The buffers dominate the polynomial coefficient growth. A multiscale choice of supports handles a fixed degree $d$, including mixed sums of old and new positions. The construction excludes every algebraic point of degree at most $d$ in an explicit positive-dimensional compact set. A degree-varying schedule gives an all-transcendental compact set with dimension zero in the Hausdorff sense and full packing dimension relative to its chosen two-symbol digit system. These are specific constructions. An arbitrary square or palindrome cap has not been shown to impose convolution buffers.

\subsection{Testing the early congruence criterion}
The independent code choices give two contrasting arithmetic families within the same bounded avoidance class. We describe these separately because neither statement assumes that a generic digit behavior holds for algebraic numbers.

\begin{proposition}\label{prop:modtest}
In every base $b\ge2$ there are uncountably many simultaneous avoiders, with square roots of length at most $2$ and palindrome lengths at most $5$, that satisfy the pairwise congruence exclusion of Corollary~\ref{cor:modular} for all fixed $p,A,c$. There is also an explicit compact class with the same caps and with
\[
 \dim_{\mathrm H}\ge (\log_b 2)/(14\cdot630)
\]
in which every point violates that exclusion at unbounded lengths and is therefore transcendental. All these constructions use the two digits $0,b-1$.
\end{proposition}

\begin{proof}
For the first statement sample the independent choices in the width-bounded code. Place the binary labels at $0,b-1$. The outcome has no six identical consecutive digits, so every tail stays at least $b^{-6}$ from each endpoint. If two length-$\ell$ block integers differ by at most $b^{(1-\varepsilon)\ell}$, their actual tails agree for the first $\lfloor\varepsilon\ell/2\rfloor$ digits for all large $\ell$. Indeed two distinct cells of this shorter depth have an inward margin of $b^{-6}$ times their width on both sides. Equal sufficiently long words in this code cannot overlap because an overlap would make a square of root length beyond $2$. Condition on all variables other than the complete tiles of the first word. The second word is fixed, which gives an exponentially small probability of an early numerical difference that small.

Next fix $p$ and a modulus $M=p^k\ge b^{\varepsilon\ell}$. Let $h$ be the lag of the compared blocks and $T=\lfloor\varepsilon\ell/4\rfloor$. For $h\ge T$ use the final $T$ digits of the later block, which are outside the earlier block. For $h<T$ use the final $T$ digits of the earlier block, which lie in the overlap. Varying complete tile choices inside the selected window changes the block difference by a change of window integer, or by $(b^h-1)$ times that change respectively. Its range in either case is smaller than $b^{2T}<M$, and distinct assignments give distinct values. Hence the chance of divisibility by $M$ is at most $2^{-\lfloor T/630\rfloor+2}$. Among starts at most $A\ell$, both bad-event probabilities are summable after the quadratic pair count. The Borel--Cantelli lemma and a countable intersection over parameters prove the first statement for almost every choice. The coding measure is atomless.

For the second statement start with the Hall--Thue fixed word for $\theta=(012,02,1)$. It begins in $0120210$. Its fourth and seventh letters give two copies of $\theta^k(0)$ after applying $\theta^k$. If $m=|\theta^k(0)|$, these disjoint source intervals begin after at least $m$ source letters and end before or at $7m$. This follows from the length recursion $(a,b,c)\mapsto(a+b+c,a+c,b)$ and the fact that $a$ is the largest. Choose successive $m$ increasing by factors at least $14$. The two copies have a common encoded length $\ell\in[605m,630m]$, and each admits $2^m$ different digit words. Fix $p\nmid b$ and take the largest $p$-power at most $2^{\lfloor m/2\rfloor}$. A residue class contains at least $2^{m/2}$ code words. A Hamming ball of radius $\lfloor m/16\rfloor$ in the $m$ choices has size at most
\[
 ((9/8)2^{3/16})^m<2^{m/2}.
\]
Choose two sufficiently distant words in one residue class and fix them in the two source intervals. The digit starts are below $8\ell$, the two words differ, and their modulus is at least $b^{(\log_b2)\ell/2520}$ for large $m$. This violates Corollary~\ref{cor:modular} on every point of the fixed class. Eventual periodicity is excluded by the square cap, so the points are transcendental. In particular these demonstrated collisions concern two words differing at a positive proportion of digit positions.

At most two length-$m$ source intervals are fixed at each stage. At least $m/2$ free source positions precede the first of its intervals, since the previous stage ends before or at $m/2$. The stage ends before or at $7m$. Thus every prefix has at least a $1/14$ free fraction, and a depth-$k$ digit cylinder has at least $\lfloor k/630\rfloor$ complete tiles. The product measure on the free choices gives the claimed dimension by the mass distribution bound.
\end{proof}

\subsection{Separating the two properties under affine maps}
The following two statements control a finite prescribed menu of slopes for occurrence and the entire rational affine family for avoidance. An affine map always has nonzero slope.

\begin{theorem}\label{thm:affine}
Fix $b\ge2$ and a finite set $\mathcal L\subset\Q\setminus\{0\}$. There are uncountably many irrational $x$ of each of the following types.
\begin{enumerate}
\item For every $\lambda\in\mathcal L$ and $t\in\Q$, $\lambda x+t$ has unbounded palindromes of each parity in base $b$. Every rational affine image of $x$ excludes arbitrary large periods $m$ of factors of length at least $wm$, for each fixed $w>1$.
\item For every $\lambda\in\mathcal L$, $t\in\Q$, and $s\ge1$, $\lambda x+t$ has unbounded square-root lengths in base $b^s$. Every rational affine image of $x$ has bounded palindrome lengths in each base $b^s$.
\end{enumerate}
Each set has an uncountable subset of transcendental numbers.
\end{theorem}

\begin{proof}
Take positive slopes first and include $1$ in the finite target menu. Write a rational affine map $f$ as $Q_f f(x)=A_f x+r_f$ with positive integer $A_f,Q_f$. Its carry
\[
 C_j=Q_f\{b^jf(x)\}-A_f\{b^jx\}-\{b^jr_f\}
\]
is an integer with $|C_j|\le A_f+Q_f+1$. For a cell of $d$ digits, \eqref{eq:digitblock} gives
\[
 Q_fD_{i,d}(f(x))\equiv A_fD_{i,d}(x)+D_{i,d}(r_f)-C_{i+d}\pmod{b^d}.
\]
Eliminate the original block from this congruence for $f$ and a map $g$ whose patterns will be avoided. The resulting congruence between the two image block integers has nonzero bounded coefficients and bounded unknown carry. With $g$ fixed the bounds are uniform over all target maps $f$ because their slope menu is finite. The rational part $r_f$ contributes only a specified block. A congruence in one unknown has at most the absolute value of its coefficient solutions modulo $b^d$. Hence, for a pair of cells, requiring the two $g$-blocks to be related by a prescribed bijection gives at most $T_g$ choices of the right target block for every left target block, for a constant independent of $d$ and of the target translations. This holds for the identity bijection as well as any prescribed reversal of groups.

To realize target windows in one real number, fix an even guard $G$ with $b^G\lambda>4$ and $b^G/\lambda>4$ for every target slope. At a stage with a closed base-$b$ input cylinder of depth $N$, leave $G$ digits before and after the target window. The image contains a complete depth-$N+G$ closed output cylinder in its interior. Appending any target word $Y$ to that output prefix yields an open cylinder whose inverse image contains a complete closed input cylinder of depth $N+2G+|Y|$ strictly inside it. This defines nested nonempty closed cylinders. The unique limit realizes all specified target windows exactly, independently of how the finite cylinders are selected.

For the first type enumerate and repeatedly use all prescribed $(\lambda,t,\text{parity})$, with lengths tending to infinity, and plant a palindrome at each window. For the second enumerate prescribed $(\lambda,t,s)$ and plant a square whose root length is a multiple of $s$. To align its start modulo $s$, precede it by an auxiliary square in the identity image. All offsets are even. The auxiliary root length can be chosen in a suitable residue class modulo $s$ so the next start is a multiple of $2s$. Guards stay fixed.

Give each reflected pair in a palindrome template, or each pair of positions at corresponding places in a templated square, an independent uniform base-$b$ digit. There are at most two positions per variable. One class per window can be fixed in advance to zero or one. We show that the unfixed digits can be chosen to avoid all unwanted long comparisons.

For the first type the comparisons are forward matches of length $k=\lfloor m/h\rfloor$ at lag $m$, for every comparison map $g_\nu$ and integer $h\ge1$. The two windows are disjoint. For the second type the comparisons are two flanks of length $k$ in base-$b$ digits, in reverse order of their $s$-digit groups. The center has width $0$ or $s$. Enumerate the comparison map, $s$, and parity jointly in $\nu$. Choose $d=d_\nu>2G$ with $T_{g_\nu}\le b^{d/2}$, and in the second type with $s\mid d$. Partition paired flanks into matching cells of width $d$ from the appropriate ends. Each match is a fixed bijection of the cell digits.

Discard cells meeting a guard, a predetermined digit, or a repeated variable in their own pair. For completeness, discard also all pairs touching a fixed initial region of length $B_\nu$ until all later planted windows have length at least $1000d$. The number of later windows touched, counted separately on each side, is at most $r=2k/(1000d)+4$. Intervening guards lose at most $4(r+4)$ pairs, and prescribed classes at most $2r$. When planted words are palindromes and the comparison is forward, internal mirror collisions require a center in the cell. A collision across a pair requires a fixed sum of the two positions, and the sums in successive cell pairs occupy disjoint ranges. When planted words are squares and the comparison is backward, an internal collision is impossible because the planted root is at least $500d$. The crosses instead require a fixed difference in disjoint ranges. In either orientation these collisions cost at most $3r$ more pairs. The early region costs at most $2B_\nu/d+4$. For $k\ge1000(B_\nu+d)$ this leaves at least $k/(2d)$ pairs with all their internal variables distinct.

Within either planted template the partners of a cell occupy at most two intervals of length at most $d$. Each interval meets at most two cells on each flank. A pair can therefore share variables with at most sixteen other pairs. Greedily retain at least $k/(34d)$ independent pairs. The probability that all of these meet the necessary congruence constraints for a genuine comparison is at most $b^{-k/68}$. This defines a finite bad event without reference to the unknown input digits.

Set $\rho=999/1000$. Impose additionally $k\ge100000+\nu+h$, with $h=1$ for the second type. Assign each event weight $\rho^k$. An event uses at most $2k$ variable classes, involving at most $4k$ positions in the whole template. For each such position, at most $2k'$ starts of a specified comparison of length $k'$ and lag qualify. In the first type the lag has at most $h'$ values for that length. Thus the neighbor weight is at most
\[
 8k\sum_{j\ge100000}(j+1)^4\rho^j<k/2000 .
\]
The inequality follows from $\rho^{1000}<1/2$ and a consecutive-term ratio below $1999/2000$ in this tail. Every event weight is at most $1/2$, and its bad probability is at most $\rho^{2k}$ since $\rho^{136}>1/2$. The product of one minus the neighboring weights is at least $\exp(-k/1000)\ge\rho^k$. The asymmetric Lov\'asz Local Lemma applies to every finite subfamily \cite[Theorem 1.2]{MT}. Compactness gives one assignment avoiding all bad events, and the nested cylinders realize it. Every actual long forbidden comparison would trigger its necessary bad event.

Taking $h$ with $1/h\le w-1$ enforces the first repetition bound. The reversed comparison enforces the second palindrome bound in its specified grouped base. A rational $x$ would violate the first type's square bound in the identity image, or the second type's palindrome bound after translating by $-x$. Thus the constructed points are irrational. For an irrational, changing sign complements the digits, which preserves all the asserted properties. Different choices of the predetermined class in each target window specify different real numbers, since that digit is fixed in the actual target image. This gives uncountability, and removing the countably many algebraic numbers gives the last assertion.
\end{proof}

\begin{remark}
Even palindromes behave differently under grouping from odd ones. For any real $x$, one even palindrome in base $b^s$ gives a square of root length $s$ in base $b$ by comparing its two central base-$b^s$ digits. If such even palindromes occur for unbounded $s$, then $\mathsf S_b(x)$ holds. An odd palindrome supplies the central relation $UVU$ with $|U|=|V|=s$, which gives a period-$2s$ factor of length $3s$. The original palindrome question does not stipulate parity.
\end{remark}

\section{Cylinder events and dimension}\label{sec:sparse-arithmetic}

The arguments below use fair digits in a reference integer base $d_0$. The reference base and the base of the patterns can differ. Write $\mathcal F_u$ for the sigma-algebra generated by the first $u$ reference digits. Let $F_{b,h}$ be the union of the closed intervals representing squares of root length $h$ and palindromes of lengths $2h$ and $2h+1$ at the first position. Write $B=b^h$. The three sets have respectively $B$, $B$, and $bB$ intervals of lengths $B^{-2}$, $B^{-2}$, and $(bB^2)^{-1}$. In particular $F_{b,h}$ has at most $(b+2)B$ interval components and measure at most $3/B$. These are the three targets of a given scale referred to below.

\begin{lemma}[Ordered cylinder events]\label{lem:ordered-cylinders}
Let $(E_i)$ be countably many clopen events on the digit space. Suppose $E_i$ depends on the first $\ell_i$ digits and that only finitely many events have closing level at most any given integer. Suppose that for an integer $0\le u_i<\ell_i$ the conditional probability of $E_i$ on each $u_i$-cylinder is at most $p_i$. Assume that some weights $0<x_i\le1/2$ satisfy
\begin{equation}\label{eq:per-level}
 \sum_{\ell_i=j}x_i\le S\quad(j\ge1),
 \qquad p_i\le x_i\exp[-2(\ell_i-u_i)S].
\end{equation}
Then a compact nonempty set of sequences avoids every $E_i$.
\end{lemma}

\begin{proof}
Order any finite subcollection by closing level, resolving ties arbitrarily. For any set of predecessors, let $P$ consist of levels at most $u_i$ and $R$ of levels strictly above it. The condition of avoiding all events in $P$ is $\mathcal F_{u_i}$-measurable. Induction on $i$ shows that the probability of $E_i$ given avoidance of any predecessor set is at most $x_i$. Indeed, after conditioning on avoidance of $P$, the numerator is bounded by $p_i$. Add the conditions in $R$ in order. By the induction hypothesis their conditional probability is at least $\prod_{j\in R}(1-x_j)$, which is at least $\exp[-2(\ell_i-u_i)S]$. This proves the claim and positivity of all conditioning denominators. Compactness gives the simultaneous conclusion. The proof is the one-sided local-lemma induction of Peres and Schlag \cite{PS}, applied with closing levels in place of a fixed index lag.
\end{proof}

\begin{lemma}[Dimension from additional cylinder constraints]\label{lem:dimension-by-coverage}
Suppose the per-level weight sum of the original constraints is at most $S_0<S$ and the second inequality in \eqref{eq:per-level} holds with $S$. The real projection of the common avoiding set has Hausdorff dimension at least $1-2S/\log d_0$.
\end{lemma}

\begin{proof}
Write $K$ for the compact projection. If $\dim_{\rm H}K<s<1-2S/\log d_0$, cover all its digit representations by cylinders of lengths $r_j$ with $\sum_j d_0^{-s r_j}<S-S_0$. To obtain such a cover, start with small real intervals and replace each by every intersecting closed reference cylinder at the comparable scale. This uses at most three cylinders of length at most $d_0$ times the interval length, including the two possible representations at grid endpoints. Compactness gives a finite subcover. Take the lengths large enough that every $d_0^{-s r_j}\le1/2$. Regard the covering cylinders as additional bad events. For one of them use $u=0$, $p=d_0^{-r_j}$ and $x=d_0^{-s r_j}$. The hypothesis on $s$ gives $p\le x e^{-2r_j S}$. With the extra events the per-level weight sums are at most $S$. Lemma~\ref{lem:ordered-cylinders} now gives a sequence avoiding the original constraints and every covering cylinder. This contradicts the cover.
\end{proof}

\section{Nonlinear rational maps and an intrinsic transcendental family}
\label{sec:nonlinear}

The integer and affine constraints in Section~\ref{sec:sparse-arithmetic} permit a common one-sided probability estimate. The same estimate survives restriction to a dyadic interval on which a nonlinear map has bounded distortion. The absolute-game argument of R5 already gives simultaneous avoidance for every algebraic-coefficient rational map. The proof below gives a direct compact construction, an effective separation for a specified algebraic input, and uniform control for prescribed moving families.

\begin{theorem}\label{thm:nonlinear}
Let $\mathcal K$ be the set of real numbers $x\in[0,1]$ with the following property. For every nonconstant $f\in\mathbb Q(X)$, the value $f(x)$ is defined and its digit expansion in every integer base has bounded square-root lengths and bounded palindrome lengths. Then $\dim_{\rm H}\mathcal K=1$, and every element of $\mathcal K$ is transcendental. For every $s<1$ the construction gives a compact subset of $\mathcal K$ of dimension at least $s$.
\end{theorem}

\begin{proof}
Use fair binary digits as the common probability space. Let $f\in\mathbb Q(X)$ be nonconstant. A point $x$ is regular for $f$ when $f$ is defined on a neighborhood of $x$ and $f'(x)\ne0$. For every binary representation of a regular point there is a closed prefix interval $J$ of depth $r$ on which $f$ has no pole, is monotone, and
\begin{equation}\label{eq:derivative-tube}
                         d\le |f'|\le 4d,\qquad d=2^k,\quad k\in\mathbb Z.
\end{equation}
Continuity permits the ratio of the maximum and minimum derivatives to be made less than $2$ by taking a sufficiently short interval. Choosing the largest dyadic power at most the minimum proves the assertion. Include in one countable list every tuple $(b,f,J,d)$ satisfying \eqref{eq:derivative-tube}, where $b\ge2$ is an integer. Different representations at a dyadic endpoint can use different intervals in this list.

Fix one tuple. Let $F_{b,h}$ be the union of the closed digit cylinders for squares of root length $h$ and palindromes of lengths $2h$ and $2h+1$. It consists of at most $(b+2)b^h$ closed arcs with total length at most $3b^{-h}$. For a shift $m$ put
\[
 q=b^m,\qquad u=\lfloor\log_2(qd)\rfloor,\qquad
 v_h=3+\lceil(2h+2)\log_2 b\rceil,\qquad \ell=u+v_h.
\]
Keep all $m$ for which $u\ge r$. The forbidden event is the union of binary cylinders of length $\ell$ with prefix $J$ whose closed real images meet
\[
       \{x\in J:q f(x)\bmod 1\in F_{b,h}\}.
\]
Its probability conditional on any binary $u$-cylinder is at most $30b^{-h}$. Indeed a $u$-cylinder either has prefix $J$ or contributes no symbolic event. On one with prefix $J$, passage to the whole covering cylinder moves $qf(x)$ by at most
$q(4d)2^{-\ell}<8\cdot2^{-v_h}\le b^{-2h-2}$. The resulting enlargement of $F_{b,h}$ has length at most $5b^{-h}$. The real interval $qf(I)$ for a binary $u$-cylinder $I$ has length between $qd|I|$ and $4qd|I|$. It meets at most $4qd|I|+2$ unit intervals, and $qd|I|\ge1$. The inverse branch has derivative bounded by $(qd)^{-1}$. The conditional probability is at most $6\cdot5b^{-h}$ as claimed. This argument uses actual real cylinders and their images. At endpoints the symbolic prefix specifies the side.

Fix $0<4\sigma<\log 2$. If the tuple has index $i$, choose an integer $H_i\ge i$ so large that
\[
       \sum_{h\ge H_i}30 b^{-h}\exp(2v_h\sigma)
                              <\sigma 2^{-i-1}.
\]
The choice is possible since $v_h\le4+(2h+2)\log_2 b$. Use the summand as the weight of the event at $(m,h)$. For each closing level $\ell$ and each $(i,h)$ at most one $m$ occurs, because $u$ increases by at least one at each step. For fixed $\ell$ only finitely many indices occur since $H_i\ge i$. The weight sum at each closing level is at most $\sigma/2$. Each event with conditional upper bound $p_i$ and weight $x_i$ satisfies $p_i\le x_i\exp(-2(\ell_i-u_i)\sigma)$. Lemma~\ref{lem:ordered-cylinders} supplies a compact set of binary sequences avoiding every event, and the covering argument of Lemma~\ref{lem:dimension-by-coverage} gives dimension at least $1-2\sigma/\log2$ for its real image.

At any point of this compact set and for any map regular there, choose a tuple using the point's binary coding. The base-$b$ tail of $f(x)$ after the finite minimum shift for that tuple avoids every target at scales $h\ge H_i$. Finite initial deletions preserve boundedness of both lengths, as shown in Section~\ref{sec:intro}. If $x$ were algebraic, its minimal polynomial would be a rational-coefficient map with a simple zero at $x$. It is regular there and has value zero. The all-zero expansion hits every target at every shift, so $x$ cannot belong to the compact set. Conversely a rational function with rational coefficients has poles or critical points only at algebraic numbers. Therefore all the maps are defined and regular at every constructed point. Letting $\sigma$ decrease proves the dimension statement.
\end{proof}

The construction gives a distance estimate for each algebraic irrational $\alpha\in(0,1)$ and each of the chosen compact sets. Use the tuple $(b,P_\alpha,J,d)$ for the minimal polynomial, with $\alpha$ in the interior of $J$. If $m_0$ is the first allowed shift and $H_i$ the scale cutoff, then every $x$ of the compact set within the interior of $J$ satisfies
\[
 \|b^{m_0}P_\alpha(x)\|>b^{-2H_i},
 \qquad
 |x-\alpha|> (4d)^{-1}b^{-m_0-2H_i}.
\]
The two end cylinders of length $b^{-2H_i}$ are among the closed square cylinders, and the mean value theorem gives the second inequality. Points outside the interior of $J$ have distance at least $\operatorname{dist}(\alpha,\partial J)$. This estimate concerns an explicitly constrained class with many polynomial conditions. For rational-coefficient maps the relevant interval tests are effective. Enumerate the charts with rational data and strict derivative bounds, decidable by polynomial sign tests on compact rational intervals. The events at each depth are finite and computable by evaluating on the two rational endpoints of each monotone cylinder. The cutoffs can be chosen by a computable geometric-series bound. Thus the separation constant is computable from the minimal polynomial, its root isolation data and the chosen construction.


\begin{theorem}[Moving families of nonlinear maps]\label{thm:moving-maps}
Fix $b\ge2$, $0\le\tau<1$ and $C_0\ge1$, and a closed binary cylinder $J$ of depth $r$. At each scale $h$ prescribe any family $\mathcal F_h$ of at most $C_0 b^{\tau h}$ real $C^1$ functions on $J$. Suppose each $f\in\mathcal F_h$ satisfies $d_f\le|f'|\le4d_f$ on $J$ for some dyadic power $d_f>0$. For all sufficiently large $H$ there is a compact $K_H\subset J$ with
\[
           \dim_{\rm H}K_H\ge 1-C_{b,\tau,C_0}b^{-(1-\tau)H}
\]
such that for every $x\in K_H$, every $h\ge H$ and every $f\in\mathcal F_h$, the expansion of $f(x)$ avoids the three targets of scale $h$ at all offsets $m\ge0$ with $b^m d_f\ge2^r$. In particular the avoidance covers every offset for maps with $d_f\ge2^r$.
\end{theorem}

\begin{proof}
Use fair binary suffixes after the prefix $J$, so the coarse and closing levels become $u-r$ and $\ell-r$. Their difference remains $v_h$, and the preceding probability bound remains $30b^{-h}$ on each coarse cylinder. A fixed map and scale contribute at most one event per closing level. Write $\lambda=(1-\tau)\log b$ and
\[
 S_H=\frac{60C_0 e^{-\lambda H}}{1-e^{-\lambda}},
 \qquad x_h=30b^{-h}e^{2v_hS_H}.
\]
Since $v_h\le4+(2h+2)\log_2 b$, the ratio to $S_H$ of the total per-level weight upper bound
\[
      30C_0\sum_{h\ge H}e^{-\lambda h+2v_hS_H}
\]
has limit at most $1/2$. This follows by summing the geometric series with rate $\lambda-4S_H\log_2 b$ and using $HS_H\to0$. For large $H$ all hypotheses of Lemmas~\ref{lem:ordered-cylinders} and \ref{lem:dimension-by-coverage} hold with slack. The real projection of the suffix space is rescaled by the length of $J$, which preserves Hausdorff dimension.
\end{proof}

Theorem~\ref{thm:moving-maps} differs from a fixed countable intersection in that one cutoff applies to the whole prescribed sequence of families. It requires no bound on the sizes of the derivatives beyond their within-map distortion, and works with functions of arbitrary complexity subject to this local hypothesis. These joint avoidance and separation statements furnish no transfer of an exact polynomial zero back to a repeat in the original digits.


\section*{Numbering correspondence}
For ease of cross-reference, the following table records the numbering of the principal statements as they appeared in an unpublished extended account of related results dated 26~November~2026.
\begin{center}
\begin{tabular}{@{}ll@{}}
\toprule
This paper & Extended account \\
\midrule
Theorem~\ref{thm:cofactor} & Theorem~4.2 \\
Corollary~\ref{cor:modular} & Corollary~4.3 \\
Proposition~\ref{prop:modtest} & Proposition~5.1 \\
Theorem~\ref{thm:affine} & Theorem~5.2 \\
Lemma~\ref{lem:ordered-cylinders} & Lemma~13.1 \\
Lemma~\ref{lem:dimension-by-coverage} & Lemma~13.2 \\
Theorem~\ref{thm:nonlinear} & Theorem~15.1 \\
Theorem~\ref{thm:moving-maps} & Theorem~15.2 \\
\bottomrule
\end{tabular}
\end{center}

\begin{thebibliography}{99}
\bibitem{AR}
B. Adamczewski and N. Rampersad,
\emph{On patterns occurring in binary algebraic numbers},
Proc. Amer. Math. Soc. 136 (2008), 3105--3109.
\href{https://doi.org/10.1090/S0002-9939-08-09319-2}{DOI}.
\href{https://adamczewski.perso.math.cnrs.fr/Overlap.pdf}{Author version}.

\bibitem{Bugeaud}
Y. Bugeaud,
\emph{Distribution Modulo One and Diophantine Approximation},
Cambridge Tracts in Mathematics, vol.~193, Cambridge University Press,
2012.
\href{https://doi.org/10.1017/CBO9781139017732}{DOI}.

\bibitem{ABind}
B. Adamczewski and Y. Bugeaud, \emph{On the independence of expansions of algebraic numbers in an integer base}, Bull. Lond. Math. Soc. 39 (2007), 283--289. \href{https://doi.org/10.1112/blms/bdl036}{DOI}. \href{https://adamczewski.perso.math.cnrs.fr/ProxNbsAlg.pdf}{Author version}.

\bibitem{BBCP}
D. H. Bailey, J. M. Borwein, R. E. Crandall, and C. Pomerance, \emph{On the binary expansions of algebraic numbers}, J. Th\'eor. Nombres Bordeaux 16 (2004), 487--518. \href{https://www.davidhbailey.com/dhbpapers/algebraic-bits.pdf}{Author version}.

\bibitem{Bilu}
Y. F. Bilu, \emph{The many faces of the subspace theorem [after
Adamczewski, Bugeaud, Corvaja, Zannier...]}, S\'eminaire Bourbaki,
expos\'e 967, Ast\'erisque 317 (2008), 1--38.
\href{https://www.numdam.org/book-part/AST_2008__317__1_0/}{Publication record}.
\href{https://arxiv.org/pdf/0907.2098#page=8}{Theorem 3.3, printed p.~8}.

\bibitem{BNY}
V. Beresnevich, E. Nesharim, and L. Yang, \emph{Winning property of badly approximable points on curves}. \href{https://arxiv.org/abs/2005.02128v3}{arXiv:2005.02128v3}.

\bibitem{DMMM}
J. Dassow, F. Manea, R. Merca\c{s}, and G. M\"uller, \emph{Inner palindromic closure}, in Developments in Language Theory, LNCS 7907, Springer, 2013. \href{https://doi.org/10.1007/978-3-642-38771-5_15}{DOI}. \href{https://oceanrep.geomar.de/id/eprint/21189/1/DasManMerMul_DLT12_CR.pdf}{Author version}.

\bibitem{EZ}
S. B. Ekhad and D. Zeilberger, \emph{There are more than $2^{n/17}$ $n$-letter ternary square-free words}, J. Integer Seq. 1 (1998), Article 98.1.9. \href{https://cs.uwaterloo.ca/journals/JIS/zeil.html}{Author text} and \href{https://cs.uwaterloo.ca/journals/JIS/zeile.html}{correction}.

\bibitem{FSU}
L. Fishman, D. S. Simmons, and M. Urba\'nski, \emph{Diophantine approximation and the geometry of limit sets in Gromov hyperbolic metric spaces}. \href{https://arxiv.org/abs/1301.5630v11}{arXiv:1301.5630v11}, Appendix C.

\bibitem{MT}
R. A. Moser and G. Tardos, \emph{A constructive proof of the general Lov\'asz Local Lemma}, J. ACM 57 (2010), Article 11. \href{https://arxiv.org/abs/0903.0544}{arXiv:0903.0544}.

\bibitem{PS}
Y. Peres and W. Schlag,
\emph{Two Erd\H{o}s problems on lacunary sequences: chromatic number and Diophantine approximation},
Bull. Lond. Math. Soc. 42 (2010), 295--300.
\href{https://doi.org/10.1112/blms/bdp126}{DOI}.
\href{https://arxiv.org/abs/0706.0223}{Author version}.


\end{thebibliography}

\end{document}
